\chapter{Πρώτες Δοκιμές}
\label{chap:protes-dokimes}

Το παρόν κεφάλαιο καταγράφει το πρώτο πειραματικό βήμα του project:
την εκπαίδευση ενός απλού convolutional δικτύου πάνω σε
log-power mel-spectrograms (mel-CNN baseline), την παρατηρηθείσα ανωμαλία
στα αποτελέσματα, και τη διαγνωστική ανάλυση που οδήγησε στον εντοπισμό
δύο σοβαρών confounders — τον codec footprint leak και την εμπλοκή φωνών
μεταξύ splits.
Η ανάλυση αυτή καθόρισε τη στρατηγική για τη συνέχεια (Phase~5+).

% ============================================================
\section{Στόχος του Baseline}
\label{sec:mel-baseline-goal}
% ============================================================

Πριν επενδυθεί χρόνος στα πολυπλοκότερα frozen audio backbones
(Wav2Vec2, WavLM, HuBERT) ή στην ανάπτυξη του multimodal pipeline,
ήταν απαραίτητο ένα \emph{γρήγορο smoke-test}: μπορεί ένα απλό μοντέλο,
που λαμβάνει ως είσοδο μόνο ηχητικά χαρακτηριστικά χαμηλού κόστους,
να διακρίνει bonafide από TTS-παραγόμενη ομιλία;

Για τον σκοπό αυτό επιλέχθηκε ένα 2D Convolutional Neural Network (CNN)
πάνω σε log-power mel-spectrograms — μία κλασική προσέγγιση στο πεδίο
του audio anti-spoofing.
Το mel-spectrogram είναι ένας συνοπτικός χρονο-συχνοτικός χάρτης
της ηχητικής ενέργειας, αποτελεσματικός για ανίχνευση διαφορών codec
και vocoder που αφήνουν χαρακτηριστικά ίχνη στο φάσμα.

Κριτήρια επιτυχίας του baseline:
\begin{itemize}
  \item Το μοντέλο να εκπαιδευτεί χωρίς σφάλματα και να συγκλίνει.
  \item Η val ROC-AUC να είναι αξιολογήσιμη (ούτε NaN ούτε αδόκιμη).
  \item Τα αποτελέσματα να παρέχουν διαγνωστική πληροφορία για το
        επόμενο βήμα σχεδίασης.
\end{itemize}

Κρίσιμη δέσμευση σε αυτή τη φάση: το \textbf{test split δεν αγγίχτηκε}.
Όλη η αξιολόγηση βασίζεται αποκλειστικά σε train και val δεδομένα.

% ============================================================
\section{Αρχιτεκτονική mel-CNN}
\label{sec:mel-cnn-arch}
% ============================================================

% ------------------------------------------------------------
\subsection{Εξαγωγή χαρακτηριστικών}
\label{subsec:mel-extraction}
% ------------------------------------------------------------

Η εξαγωγή mel-spectrograms υλοποιείται στο \codeid{src/features/extract\_mel.py}
μέσω της κλάσης \codeid{MelExtractor}.
Κάθε δείγμα ήχου φορτώνεται μέσω \codeid{src/features/audio\_io.py}
(deterministic 4{,}0~s παράθυρο στα 16~kHz mono),
στη συνέχεια περνάει από
\codeid{torchaudio.transforms.MelSpectrogram} και \codeid{AmplitudeToDB},
και αποθηκεύεται ως \codeid{.npy} αρχείο σχήματος \((n\_\text{mels},\ T)\).

Οι παράμετροι εξαγωγής προέρχονται από το \codeid{config/default.yaml}:
\begin{itemize}
  \item \codeid{sample\_rate}: 16\,000~Hz
  \item \codeid{n\_fft}: 400, \codeid{win\_length}: 400, \codeid{hop\_length}: 160
  \item \codeid{n\_mels}: 64
  \item \codeid{seconds}: 4{,}0
\end{itemize}

Με αυτές τις παραμέτρους, κάθε δείγμα παράγει ένα spectrogram
\((64,\ 401)\) — 64 mel bins και 401 χρονικά frames.
Τα \codeid{.npy} αρχεία αποθηκεύονται στον κατάλογο
\codeid{data/features/audio\_mel/}, ένα ανά \codeid{sample\_id}.
Το raw WAV \textbf{δεν} εισέρχεται ποτέ στον training loop — μόνο τα
προεξαγμένα \codeid{.npy} features.

% ------------------------------------------------------------
\subsection{Αρχιτεκτονική δικτύου}
\label{subsec:mel-cnn-model}
% ------------------------------------------------------------

Το \codeid{MelCNN} είναι ένα compact 1-channel 2D CNN με τέσσερα
blocks \codeid{Conv2d → BatchNorm2d → ReLU → MaxPool2d},
με κανάλια (16, 32, 64, 128), ακολουθούμενα από
\codeid{AdaptiveAvgPool2d(1)} → flatten → \codeid{Dropout(0.3)} →
\codeid{Linear(128 → 1)} logit.

Το σύνολο των εκπαιδεύσιμων παραμέτρων ανέρχεται σε \textbf{97\,521},
εντός του project budget των 2\,000\,000 παραμέτρων.

% ------------------------------------------------------------
\subsection{Normalization και training}
\label{subsec:mel-training}
% ------------------------------------------------------------

Τα normalization statistics (mean/std ανά mel bin) υπολογίστηκαν
\textbf{αποκλειστικά από το train split} μέσω της \codeid{fit\_mel\_norm\_stats}.
Τα val δεδομένα χρησιμοποιήθηκαν μόνο για αξιολόγηση — ποτέ για
υπολογισμό στατιστικών.

Hyperparameters εκπαίδευσης από \codeid{config/default.yaml}:
\begin{itemize}
  \item Optimizer: AdamW, \codeid{lr}=1e-4, \codeid{weight\_decay}=1e-2
  \item Loss: \codeid{BCEWithLogitsLoss}
  \item \codeid{batch\_size}: 32, \codeid{max\_epochs}: 50
  \item Early stop: patience 7 εποχές βάσει val ROC-AUC
  \item Seed: 42
\end{itemize}

Το τελικό checkpoint (\codeid{models/checkpoints/best\_mel\_cnn.pt})
αποθηκεύει αυτοπεριγραφικά τα weights, τη mel config, τα normalization
statistics και τα val metrics, ώστε η ανακατασκευή του να μην απαιτεί
επαναρύθμιση.

% ============================================================
\section{Τα Πρώτα Αποτελέσματα ήταν Ύποπτα}
\label{sec:results-suspicious}
% ============================================================

Το μοντέλο εκπαιδεύτηκε με early stopping και σταμάτησε στην
\textbf{εποχή 17} (patience εξαντλήθηκε μετά από 7 εποχές χωρίς βελτίωση).
Οι επιδόσεις στο val split:

\begin{center}
\begin{tabular}{lr}
\hline
Μετρική & Τιμή \\
\hline
Val ROC-AUC & 1{,}0000 \\
Val EER & 0{,}0000 \\
Val F1 (@ EER threshold) & 1{,}0000 \\
Val Precision & 1{,}0000 \\
Val Recall & 1{,}0000 \\
\hline
\end{tabular}
\end{center}

Η val ROC-AUC έφτασε 1{,}0 ήδη από την εποχή 10.
Η per-provider recall ήταν επίσης 1{,}0 και για τους δύο TTS providers
(\engineName{ElevenLabs}, \engineName{Google~TTS}), καθώς και για τις
bonafide (\codeid{original}) εγγραφές.

Αυτά τα αποτελέσματα \textbf{δεν ερμηνεύθηκαν ως θρίαμβος} —
αντίθετα, ερμηνεύθηκαν ως \emph{διαγνωστικό συναγερμό}.

Ένα compact CNN των 97\,521 παραμέτρων, εκπαιδευμένο για μόλις 17 εποχές
σε dataset 1\,518 δειγμάτων, δεν αναμενόταν να φτάσει σε τέλεια διαχωρισμό.
Η μόνη εξήγηση ήταν ότι το μοντέλο \textbf{δεν έμαθε deepfake detection} —
έμαθε να αναγνωρίζει έναν πολύ πιο εύκολο διαχωρισμό που ήταν
λανθάνοντα ενσωματωμένος στα δεδομένα.

% ============================================================
\section{Δύο Confounders που Αποκαλύφθηκαν}
\label{sec:confounders}
% ============================================================

Η έρευνα που ακολούθησε εντόπισε δύο ανεξάρτητους confounders
που μπορούσαν να εξηγήσουν την τέλεια val ROC-AUC.

% ------------------------------------------------------------
\subsection{Codec Footprint Leak}
\label{subsec:codec-leak}
% ------------------------------------------------------------

Ο πρώτος — και σοβαρότερος — confounder ήταν πλήρης συσχέτιση
μορφότυπου αρχείου με ετικέτα κλάσης:

\begin{itemize}
  \item \textbf{Bonafide} δείγματα: αποθηκευμένα ως clean 16~kHz PCM WAV —
        χωρίς lossy compression.
  \item \textbf{TTS spoof} δείγματα: παραδοτέα ως lossy MP3
        (\engineName{ElevenLabs}: 44{,}1~kHz / 128~kbps·
         \engineName{Google~TTS}: 24~kHz / 64~kbps).
\end{itemize}

Αυτή η 100\% συσχέτιση \codeid{format ↔ label} σημαίνει ότι
\textbf{οποιοδήποτε μοντέλο που δέχεται mel-spectrograms ως είσοδο
μπορεί να λειτουργήσει ως codec detector} αντί για deepfake detector.
Τα lossy MP3 artifacts — bandwidth ceiling, quantization noise,
aliasing στις υψηλές συχνότητες — είναι ορατά στο mel-spectrogram και
διαφέρουν δραματικά από το clean WAV bonafide.

Το mel-CNN δεν χρειάστηκε να μάθει τίποτα σχετικό με την αυθεντικότητα
ομιλίας: αρκούσε να ανιχνεύει codec footprint.
Αυτός ο shortcut εξηγεί πλήρως την τέλεια val ROC-AUC.

% ------------------------------------------------------------
\subsection{Voice Straddling Splits}
\label{subsec:voice-straddling}
% ------------------------------------------------------------

Ο δεύτερος confounder εντοπίστηκε στη δομή των splits.
Ακόμη και αν ο codec leak επιδιορθωνόταν, η ίδια TTS φωνή
(π.χ. ένα συγκεκριμένο \codeid{voice\_id} του \engineName{ElevenLabs})
μπορούσε να εμφανίζεται ταυτόχρονα σε train, val και test.

Αυτό σημαίνει ότι το μοντέλο μπορούσε να μάθει \emph{φωνο-ειδικά}
χαρακτηριστικά στο train και να τα αναγνωρίσει στο val —
όχι γιατί κατανοεί τη διαδικασία TTS-σύνθεσης,
αλλά γιατί \textbf{έχει ακούσει αυτή τη φωνή κατά την εκπαίδευση}.

Αυτό το πρόβλημα είναι γνωστό ως \emph{voice straddling across splits}
και αποτελεί ένα βασικό σφάλμα στον σχεδιασμό συνόλων δεδομένων
αντι-spoofing: το evaluation protocol πρέπει να εξασφαλίζει ότι κάθε
\codeid{(provider, voice\_id)} βρίσκεται αποκλειστικά σε ένα split.

% ============================================================
\section{Συμπεράσματα από τη Διάγνωση}
\label{sec:phase4-diagnosis}
% ============================================================

Τα δύο αυτά ευρήματα κατέστησαν σαφές ότι τα αποτελέσματα του
mel-CNN baseline \textbf{δεν έχουν διαγνωστική αξία για deepfake detection}:
αντικατοπτρίζουν codec identification και πιθανή φωνο-αναγνώριση,
όχι ανίχνευση συνθετικής ομιλίας.

Καθορίστηκαν τρία διορθωτικά βήματα για τη συνέχεια:

\paragraph{1. Codec matching.}
Κάθε bonafide WAV περνάει μέσω ενός τυχαία επιλεγμένου MP3
round-trip (κωδικοποίηση σε MP3 → αποκωδικοποίηση πίσω σε 16~kHz mono WAV),
χρησιμοποιώντας specs που δειγματοληπτούνται από την κατανομή των
TTS-spoof αρχείων.
Έτσι το codec history καθίσταται ανεξάρτητο από την ετικέτα.
Τα TTS spoofs που παράγονται εγγενώς ως clean WAV
(π.χ. \engineName{OpenAI~TTS}) υπόκεινται επίσης στο ίδιο round-trip,
ώστε να μην ανατρέψουν τον διαχωρισμό σε αντίθετη κατεύθυνση.

\paragraph{2. Voice-disjoint split assignment.}
Κάθε \codeid{(provider, voice\_id)} τοποθετείται αποκλειστικά σε ένα split.
Η απόφαση split ταξιδεύει από το audio manifest στα visual και fusion
manifests μέσω \codeid{src/data/apply\_voice\_split.py},
ώστε όλες οι views να συμφωνούν πλήρως.

\paragraph{3. Anti-leakage audit ως υποχρεωτικό βήμα pipeline.}
Η εμπειρία από τον mel-CNN baseline αξιοποιήθηκε για να κωδικοποιηθεί
ένα εκτεταμένο anti-leakage audit ως τυπικό βήμα του pipeline,
που εκτελείται πριν από κάθε εκπαίδευση μέσω
\codeid{python -m src.data.feature\_store --validate}.

% ============================================================
\section{Από raw mel σε Frozen Embeddings}
\label{sec:mel-to-embeddings}
% ============================================================

Παράλληλα με τη διόρθωση των confounders, αναθεωρήθηκε και η αρχιτεκτονική
επιλογή: η μετάβαση από raw mel-spectrograms σε \emph{frozen pretrained embeddings}.

Η λογική της αλλαγής στηρίζεται σε τρεις λόγους:

\begin{enumerate}
  \item \textbf{Κατώτατο param budget.}
        Τα mel-CNN features είναι time-frequency textures·
        για την εκμάθησή τους απαιτείται αρκετά βαθύ CNN.
        Αντίθετα, οι frozen embeddings Wav2Vec2 / WavLM / HuBERT
        περιγράφουν ήδη πλούσιες φωνητικές αναπαραστάσεις
        σε 768-d διανύσματα ανά χρονικό frame.
        Ένας απλός MLP head \((768 \to 256 \to 128)\) αρκεί,
        διατηρώντας τον αριθμό εκπαιδεύσιμων παραμέτρων χαμηλό.

  \item \textbf{Αναπαραγωγιμότητα.}
        Τα frozen backbones παράγουν ντετερμινιστικά τα ίδια embeddings
        κάθε φορά.
        Η εξαγωγή γίνεται \emph{μία φορά} και αποθηκεύεται ως \codeid{.npy}·
        η εκπαίδευση διαβάζει πάντα από αυτά τα αρχεία,
        χωρίς ποτέ να αγγίζει το raw audio στον training loop.

  \item \textbf{Transfer από μεγαλύτερα pretrained μοντέλα.}
        Τα \codeid{facebook/wav2vec2-base-960h},
        \codeid{microsoft/wavlm-base} και
        \codeid{facebook/hubert-base-ls960} έχουν εκπαιδευτεί
        σε εκατοντάδες έως χιλιάδες ώρες ομιλίας.
        Η γνώση αυτή μεταφέρεται δωρεάν στο task μας —
        χωρίς fine-tuning, χωρίς κίνδυνο catastrophic forgetting,
        και με το πλήρες ιστορικό pretraining να παραμένει αμετάβλητο.
\end{enumerate}

Τα backbones παραμένουν \textbf{πλήρως παγωμένα} (frozen) καθ' όλη τη
διάρκεια της εκπαίδευσης.
Ο αριθμός εκπαιδεύσιμων παραμέτρων για το audio head περιορίζεται
σε \todoMetric{περίπου 328\,000} — κατά πολύ κάτω του project budget.
Ο ίδιος αρχιτεκτονικός περιορισμός (frozen backbone + ελαφριά head)
εφαρμόζεται και στους υπόλοιπους modalities (visual BiGRU, fusion MLP).
